% ============================================================
% style/page-style.tex
%
% Базовое оформление страницы.
%
% Здесь задаются:
%   - цвет фона страницы;
%   - основной цвет текста;
%   - стиль нумерации страниц;
%   - параметры абзацев;
%   - оформление подписей;
%   - цвет линий в таблицах.
%
% Важно:
%   этот файл должен подключаться ПОСЛЕ:
%     style/colors.tex
%
% Потому что здесь используются смысловые цвета:
%   rbg, rtext, rframe.
% ============================================================


% ------------------------------------------------------------
% 1. Фон страницы и основной цвет текста
% ------------------------------------------------------------
%
% rbg    — основной фон документа;
% rtext  — основной цвет текста.
%
% В colors.tex сейчас:
%   rbg   = pag  = #293133;
%   rtext = white.

\pagecolor{rbg}
\color{rtext}


% ------------------------------------------------------------
% 2. Стиль страниц
% ------------------------------------------------------------
%
% Без fancyhdr, декоративных линий и рамок страницы.
%
% plain:
%   - номер страницы снизу по центру;
%   - без верхнего колонтитула.

\pagestyle{plain}


% ------------------------------------------------------------
% 3. Абзацы
% ------------------------------------------------------------
%
% Стиль близок к стандартному report/article:
%   - первая строка абзаца имеет отступ;
%   - дополнительного вертикального расстояния между абзацами нет.

\setlength{\parindent}{15pt}
\setlength{\parskip}{0pt}


% ------------------------------------------------------------
% 4. Переносы и переполнение строк
% ------------------------------------------------------------
%
% \tolerance управляет тем, насколько TeX готов растягивать пробелы,
% чтобы избежать плохих переносов.
%
% \emergencystretch даёт TeX дополнительный запас растяжения строки.
% Это помогает уменьшить количество overfull \hbox.

\tolerance=1000
\emergencystretch=2em


% ------------------------------------------------------------
% 5. Подписи к рисункам и таблицам
% ------------------------------------------------------------
%
% labelformat=empty:
%   убирает стандартные подписи вида "Рис. 1", "Таблица 1".
%
% labelsep=none:
%   убирает разделитель после номера подписи.
%
% justification=centering:
%   центрирует подпись.
%
% singlelinecheck=true:
%   одиночная строка подписи дополнительно центрируется.
%
% textfont={color=rtext}:
%   делает текст подписи основным цветом документа.

\captionsetup{
    labelformat=empty,
    labelsep=none,
    justification=centering,
    singlelinecheck=true,
    font=normalsize,
    textfont={color=rtext}
}


% ------------------------------------------------------------
% 6. Линии в таблицах
% ------------------------------------------------------------
%
% По умолчанию линии таблиц делаем цветом декоративной рамки.
% Сейчас rframe = white.

\arrayrulecolor{rframe}